% sections/conclusion.tex

\section{图片及表格示例}
% ====================================================
% 本节展示论文中使用的图表示例和格式规范
% 所有图表都采用中英文双语标题，符合学报要求
% ====================================================

% ==================== 表格示例部分 ====================
\subsection{表格示例}
% 表格使用说明：所有表格采用三线表格式，表头使用英文
% 表格标题使用中英文双语，表格内容使用英文

% 表格引用说明：使用 \ref{} 命令引用表格编号
实验环境配置详见表 \ref{tab:experiment_environment}，所有测试均在统一硬件和软件环境下进行，确保结果的可比性和可重复性。

\begin{table}[htbp]
	% 表格位置参数说明：
	% [htbp] 让 LaTeX 自动选择最佳位置
	% h=此处, t=顶部, b=底部, p=单独页
	
	\centering % 表格居中
	\tabbicaption{实验环境配置详情} % 中文标题
	{Experimental environment configuration details} % 英文标题
	\label{tab:experiment_environment} % 标签用于交叉引用
	
	\begin{tabular}{ll} % 两列都左对齐
		\toprule % 上横线（三线表顶部）
		\textbf{Configuration Item} & \textbf{Detailed Parameters} \\ % 英文表头
		\midrule % 中间横线（分隔表头和内容）
		Processor & Intel Core i7-11800H @ 2.30GHz \\
		Memory & 32GB DDR4 3200MHz \\
		Operating System & Ubuntu 20.04.3 LTS \\
		Programming Language & Python 3.8.10 \\
		Math Library & NumPy 1.21.2, SymPy 1.9 \\
		\bottomrule % 底部横线（三线表底部）
	\end{tabular}
	% 注意：需要 booktabs 宏包支持三线表
\end{table}

根据实验数据，我们对四种算法进行了全面测试。测试环境配置如下：
\begin{itemize}
	\item 处理器：Intel i7-11800H @ 2.30GHz
	\item 内存：32GB DDR4 3200MHz
	\item 操作系统：Ubuntu 20.04 LTS
\end{itemize}

% ==================== 性能对比表格示例 ====================
\begin{table}[H] % 使用H强制位置
	\centering
	\tabbicaption{算法性能对比结果}
	{Performance comparison results of different algorithms}
	\label{tab:performance_comparison}
	
	% 使用resizebox自动缩放（保持原始比例）
	\resizebox{\textwidth}{!}{ % 开始缩放
		\begin{tabular}{l *{2}{S[table-format=2.1]} *{3}{S[table-format=3.0]}} % 5列数据
			\toprule
			\multicolumn{1}{c}{\textbf{算法}} & 
			\multicolumn{1}{c}{\textbf{时间(ms)}} & 
			\multicolumn{1}{c}{\textbf{加速比}} & 
			\multicolumn{1}{c}{\textbf{内存(MB)}} & 
			\multicolumn{1}{c}{\textbf{吞吐量(Mbps)}} & 
			\multicolumn{1}{c}{\textbf{准确率(\%)}} \\
			\midrule
			传统Montgomery & 12.5 & 1.0 & 256 & 120 & 99.8 \\
			改进算法 & 8.7 & 1.4 & 198 & 180 & 99.9 \\
			优化版本 & 6.2 & 2.0 & 215 & 210 & 99.9 \\
			混合方案 & 9.5 & 1.3 & 230 & 165 & 99.7 \\
			\bottomrule
		\end{tabular}
	} % 结束缩放
	
	% 添加表格注释（使用threeparttable包）
	\begin{tablenotes}
		\small
		\item[1] 测试环境：Intel i7-11800H @ 2.30GHz, 32GB DDR4
		\item[2] 数据为10次测试平均值
		\item[3] 吞吐量测试使用1024-bit模数
	\end{tablenotes}
\end{table}

从表\ref{tab:performance_comparison}可以得出以下结论：
\begin{enumerate}
	\item 优化版本在时间效率上表现最佳，比传统算法快2倍
	\item 改进算法在内存占用上最优，减少约22.6%
	\item 混合方案在吞吐量和内存之间取得较好平衡
\end{enumerate}

% ==================== 图片示例部分 ====================
\subsection{图片示例}
% 图片使用说明：所有图片采用矢量格式（PDF/EPS）
% 图片标题使用中英文双语，确保清晰度

% 图片引用说明：使用 \ref{} 命令引用图片编号
图 \ref{fig:performance_comparison} 展示了不同算法在相同测试数据集上的性能对比结果。从图中可以看出，本文提出的算法在时间效率上显著优于传统方法。

\begin{figure}[htbp]
	\centering
	\figbicaption{各算法性能对比分析图} % 中文标题
	{Performance comparison analysis of different algorithms} % 英文标题
	\label{fig:performance_comparison} % 标签
	\includegraphics[width=0.85\textwidth]{figures/figure2.eps} % 矢量图格式
	% 图片尺寸说明：使用相对尺寸（0.85\textwidth）确保自适应
\end{figure}

% 第二个图片示例：算法执行流程图
图 \ref{fig:algorithm_flow} 展示了本文算法的完整执行流程，包括输入处理、核心计算和结果输出三个阶段。

\begin{figure}[H] % 使用H强制当前位置
	\centering
	\figbicaption{改进Montgomery算法的执行流程图} % 中文标题
	{Execution flow chart of improved Montgomery algorithm} % 英文标题
	\label{fig:algorithm_flow} % 标签
	\includegraphics[width=0.85\textwidth]{figures/ECUST-logo.png} % PDF矢量格式
	\end{figure}
	
% ==================== 交叉引用示例 ====================
% 交叉引用说明：使用 \ref{} 引用图表编号
如图 \ref{fig:algorithm_flow} 所示，算法通过三个主要阶段实现了有限域上的高效模乘运算。结合表 \ref{tab:experiment_environment} 中的实验环境数据和表 \ref{tab:performance_comparison} 的性能结果，可以全面评估算法的有效性。

% ==================== 子图表示例 ====================
% 子图片示例：使用 subcaption 宏包支持多个子图
\begin{figure}[H] % 使用H强制当前位置
	\centering
		% 主图标题
	\figbicaption{算法复杂度对比分析} % 主图中文标题
	{Comparative analysis of algorithm complexity} % 主图英文标题
	\label{fig:complexity_analysis} % 主图标签
	
	% 第一个子图：时间复杂度分析
	\begin{subfigure}{0.45\textwidth} % 子图宽度
		\centering
		\figbicaption{时间复杂度分析} % 子图中文标题
		{Time complexity analysis} % 子图英文标题
		\label{fig:time_complexity} % 子图标签
		\includegraphics[width=0.85\linewidth]{figures/ECUST-logo.png}
	\end{subfigure}% 子图中间不要随便添加空格换行，会导致空间不够第二子图自动换行
	\hfill % 填充间距
	% 第二个子图：空间复杂度分析
	\begin{subfigure}{0.45\textwidth}
		\centering
		\figbicaption{空间复杂度分析} % 子图中文标题
		{Space complexity analysis} % 子图英文标题
		\label{fig:space_complexity} % 子图标签
		\includegraphics[width=0.85\linewidth]{figures/ECUST-logo.png}
	\end{subfigure}

\end{figure}



时间复杂度和空间复杂度的详细分析分别见图 \ref{fig:time_complexity} 和图 \ref{fig:space_complexity}，整体复杂度对比见图 \ref{fig:complexity_analysis}。
